\documentclass[11pt]{article}

\usepackage[T1]{fontenc}
\usepackage[utf8]{inputenc}
\usepackage{lmodern}
\usepackage{microtype}
\usepackage[margin=1in]{geometry}
\usepackage{booktabs}
\usepackage{tabularx}
\usepackage{amsmath}
\usepackage{enumitem}
\usepackage[dvipsnames]{xcolor}
\usepackage{hyperref}

\definecolor{RoverBlue}{HTML}{1F4E79}
\definecolor{HazardOrange}{HTML}{C55A11}
\definecolor{SoftGray}{HTML}{F2F2F2}

\hypersetup{
  colorlinks=true,
  linkcolor=RoverBlue,
  urlcolor=RoverBlue,
  pdftitle={Hazard Inspection Rover: Six-Module Course Outline},
  pdfauthor={hazard_inspection_rover contributors}
}

\setlist[itemize]{leftmargin=1.5em, itemsep=0.25em, topsep=0.35em}
\setlist[enumerate]{leftmargin=1.7em, itemsep=0.25em, topsep=0.35em}
\renewcommand{\arraystretch}{1.2}

\newcommand{\projectname}{\texttt{hazard\_inspection\_rover}}
\newcommand{\ros}{ROS~2 Humble}
\newcommand{\gazebo}{Gazebo Harmonic}
\newcommand{\keypoint}[1]{%
  \begin{center}
    \fcolorbox{HazardOrange}{SoftGray}{%
      \parbox{0.9\textwidth}{\textbf{Key point:} #1}%
    }
  \end{center}
}

\title{%
  \textcolor{RoverBlue}{\textbf{Hazard Inspection Rover}}\\[0.35em]
  \large Six-Module ROS 2 Course Outline
}
\author{Beginner-to-Intermediate Robotics Portfolio Course}
\date{Target environment: Ubuntu 22.04.5 LTS, WSL2, Python 3.10, ROS 2 Humble}

\begin{document}

\maketitle

\begin{abstract}
This guided course develops a small differential-drive inspection-rover
simulation as a credible robotics software portfolio project. It is designed
for a mechanical engineer who understands physical systems, modeling, testing,
and failure analysis but has limited practical ROS~2 experience. The project
emphasizes system integration, explicit interfaces, simple feedback control,
defensive behavior, and evidence-based validation. It does not claim SLAM,
artificial intelligence, machine learning, optimal control, full autonomy, or
field readiness.
\end{abstract}

\tableofcontents
\newpage

\section{Course purpose and engineering context}

The completed \projectname{} project simulates a small differential-drive
mobile robot performing a predefined inspection route in a simplified
industrial environment. The rover publishes odometry, laser scans, joint
states, transforms, mission state, and a mathematically simulated hazard
reading. A deliberately simple proportional waypoint controller moves the
robot between inspection stations. If the hazard level becomes excessive, or
if required sensor data becomes stale or invalid, the controller commands a
stop and enters a defined retreat or abort behavior. Mission data are recorded
to CSV for later inspection.

For a mechanical engineer, the important transition is from analyzing an
individual physical component to reasoning about a distributed robotic system.
The course connects familiar engineering habits---coordinate systems,
free-body reasoning, parameter selection, tolerances, failure modes, test
evidence, and traceability---to ROS~2 software concepts.

\subsection{Technical scope}

The target software environment is:

\begin{itemize}
  \item Ubuntu 22.04.5 LTS (Jammy) under WSL2;
  \item Python 3.10 through \texttt{python3};
  \item \ros{};
  \item \gazebo{}, using \texttt{ros\_gz} integration rather than Gazebo
        Classic;
  \item RViz2 and VS Code connected to WSL;
  \item no physical hardware, proprietary software, ROS~1, or required Docker
        environment.
\end{itemize}

\subsection{Learning philosophy}

The course follows the progression

\begin{center}
\textbf{communicate $\rightarrow$ model $\rightarrow$ simulate
$\rightarrow$ decide $\rightarrow$ integrate $\rightarrow$ validate}.
\end{center}

Each module adds one major source of complexity while preserving the student's
ability to inspect the system. A module ends with a verification gate: the
student predicts behavior, runs the system, collects evidence, explains the
result, and makes a small modification. The next module should begin only after
that gate is passed.

\section{Selected software architecture}

The project uses five packages with explicit responsibilities.

\begin{table}[htbp]
\centering
\small
\begin{tabularx}{\textwidth}{@{}l l X@{}}
\toprule
\textbf{Package} & \textbf{Build type} & \textbf{Responsibility} \\
\midrule
\texttt{hazard\_inspection\_description}
  & \texttt{ament\_cmake}
  & Xacro robot model, primitive geometry, inertial properties, and RViz2
    configuration. \\
\texttt{hazard\_inspection\_sim}
  & \texttt{ament\_python}
  & Gazebo world and simulation assets, spawning support, and mathematical
    hazard-sensor node. \\
\texttt{hazard\_inspection\_control}
  & \texttt{ament\_python}
  & Waypoint controller, finite-state mission logic, mission services, CSV
    logger, and testable control calculations. \\
\texttt{hazard\_inspection\_interfaces}
  & \texttt{ament\_cmake}
  & A small, strongly typed mission-state message. Standard services are used
    where they are sufficient. \\
\texttt{hazard\_inspection\_bringup}
  & \texttt{ament\_python}
  & Main launch file, integrated parameter files, and ROS--Gazebo bridge
    configuration. \\
\bottomrule
\end{tabularx}
\caption{Target ROS 2 package responsibilities.}
\end{table}

This structure is intentionally more disciplined than placing every file in a
single package. It teaches dependency ownership and makes it possible to test
control logic without loading robot descriptions or launching the simulator.
The additional packaging overhead is small enough for a beginner and produces
a clearer portfolio repository.

\section{Why the course uses this order}

The ordering is dependency-driven and risk-driven.

\begin{enumerate}
  \item \textbf{Communication comes first.} A student must be able to build,
        source, run, and inspect ROS~2 nodes before debugging a robot model or
        simulator. Otherwise, a simple workspace error can be mistaken for a
        physics or control problem.
  \item \textbf{The robot model precedes simulation.} Frames, joint axes,
        dimensions, collision shapes, and inertial properties define what the
        simulator will act on. RViz2 provides a lower-complexity environment in
        which to verify that structure before physics is introduced.
  \item \textbf{Simulation precedes mission behavior.} Odometry and laser data
        are the controller's inputs. Establishing their topic names, units,
        frames, signs, and update rates first prevents control debugging from
        becoming mixed with bridge debugging.
  \item \textbf{Hazard sensing and state logic are isolated before motion is
        integrated.} A safety response should be testable as a decision system,
        not only by watching a moving robot. This allows deterministic testing
        of thresholds, invalid inputs, and services.
  \item \textbf{Full motion is integrated only after interfaces are stable.}
        The waypoint controller then connects known inputs and outputs. Problems
        can be localized to sensing, decision logic, actuation, or configuration.
  \item \textbf{Portfolio work follows technical validation.} Screenshots and
        claims should describe demonstrated behavior. Tests and a validation
        matrix are completed before the final README, video, resume bullets, and
        statement-of-purpose descriptions are treated as evidence.
\end{enumerate}

This order also mirrors a conservative engineering development process:
component checks are followed by subsystem checks, integration, fault
injection, and documented acceptance evidence.

\section{Course-at-a-glance}

\begin{table}[htbp]
\centering
\small
\begin{tabularx}{\textwidth}{@{}c l c X@{}}
\toprule
\textbf{Module} & \textbf{Theme} & \textbf{Time} & \textbf{Primary evidence} \\
\midrule
1 & ROS 2 fundamentals & 2--3 h
  & Two discoverable nodes exchange a typed message at the predicted rate. \\
2 & Robot model and TF & 3--5 h
  & A dimensionally consistent rover and correct frame tree appear in RViz2. \\
3 & Gazebo simulation & 4--5 h
  & The rover spawns, accepts velocity commands, and publishes odometry and
    laser scans. \\
4 & Hazard and mission logic & 3--5 h
  & Hazard values are reproducible; services and safety-related state
    transitions behave as specified. \\
5 & Mission integration & 4--5 h
  & One launch command produces waypoint motion, inspection stops, hazard
    retreat, RViz2 data, and CSV output. \\
6 & Validation and portfolio & 3--5 h
  & Tests, validation records, documentation, and a factual demonstration form
    a reproducible portfolio artifact. \\
\bottomrule
\end{tabularx}
\caption{Six-module learning sequence.}
\end{table}

\newpage
\section{Module 1: ROS 2 workspace and communication fundamentals}

\textbf{Estimated time:} 2--3 hours

\subsection{Purpose}

Module 1 creates the smallest useful ROS~2 system in the course: one Python node
publishes a numbered inspection-status string, and a second node subscribes and
logs it. The example uses a separate training topic so it cannot be confused
with the final strongly typed mission-state interface.

\subsection{Why it comes first}

Every later feature depends on the same fundamental workflow: package metadata
must be correct, \texttt{colcon} must install an executable, each terminal must
source the correct environment, and the ROS graph must connect matching topic
names and message types. Learning this with two small nodes produces a clean
diagnostic baseline.

\subsection{Learning objectives}

The student will be able to:

\begin{itemize}
  \item distinguish a workspace, package, Python module, executable, node,
        topic, and message instance;
  \item build one package and source the resulting overlay;
  \item explain how \texttt{package.xml}, \texttt{setup.py}, the resource
        marker, and console entry points support discovery;
  \item create a timer-driven publisher and callback-driven subscriber;
  \item inspect nodes, topics, types, endpoints, parameters, and message rates
        with ROS~2 command-line tools;
  \item validate configuration instead of silently accepting invalid values.
\end{itemize}

\subsection{Core concepts}

ROS graph discovery, publishers, subscribers, typed messages, callback
execution, QoS queue depth, parameters, launch files, package installation,
base environments, and workspace overlays.

\subsection{Principal deliverables}

\begin{itemize}
  \item \texttt{hazard\_inspection\_control} Python package;
  \item \texttt{inspection\_status\_publisher} executable;
  \item \texttt{inspection\_status\_monitor} executable;
  \item \texttt{/training/inspection\_status} String topic;
  \item introductory YAML parameter file and two-node launch file;
  \item recorded command-line evidence and a small parameter/code exercise.
\end{itemize}

\keypoint{ROS 2 components do not communicate because their Python files know
about one another. They communicate through discoverable, named, typed
interfaces. Package installation and environment sourcing make those
components available to ROS 2 tools.}

\subsection{Verification gate}

The student must show both nodes in \texttt{ros2 node list}, inspect the topic's
type and endpoints, echo a message, measure approximately 1~Hz at the default
setting, modify the configured rate, and demonstrate that a non-positive timer
period is rejected.

\subsection{Mechanical-engineering connection}

The topic contract is analogous to an instrument interface: signal identity,
units or representation, update rate, and expected consumers must agree. A
running process is not evidence that the measurement chain is connected.

\subsection{Suggested commits}

\begin{itemize}
  \item \texttt{chore: initialize ROS 2 workspace}
  \item \texttt{feat: add basic ROS 2 communication nodes}
\end{itemize}

\section{Module 2: Robot model, coordinate frames, and RViz2}

\textbf{Estimated time:} 3--5 hours

\subsection{Purpose}

Module 2 describes the rover using Xacro and primitive geometry. The model
includes visual, collision, and inertial definitions together with
\texttt{base\_footprint}, \texttt{base\_link}, wheel, caster,
\texttt{laser\_frame}, and \texttt{sensor\_frame} links. RViz2 is used to
inspect the model and transform tree without simulation physics.

\subsection{Why it follows Module 1}

Robot-state publishing is itself a ROS~2 data flow. The student now applies the
communication and inspection workflow from Module 1 to a structured mechanical
model. Separating RViz verification from Gazebo keeps geometry and transform
errors distinct from contact and dynamics errors.

\subsection{Learning objectives}

The student will be able to:

\begin{itemize}
  \item explain the difference between a link and a joint;
  \item use Xacro properties and macros to avoid repeated dimensions;
  \item define realistic simple masses, centers of mass, and inertia tensors;
  \item distinguish visual geometry from collision geometry;
  \item define joint origins and axes using ROS coordinate conventions;
  \item explain every frame in the rover's TF tree;
  \item use \texttt{robot\_state\_publisher}, RViz2, \texttt{tf2\_echo}, and
        frame-graph tools to verify the model.
\end{itemize}

\subsection{Core concepts}

URDF, Xacro, links, joints, origins, roll--pitch--yaw, inertial properties,
collision geometry, visual geometry, \texttt{/tf}, \texttt{/tf\_static}, and
the distinction between robot state and world pose.

\subsection{Principal deliverables}

\begin{itemize}
  \item \texttt{hazard\_inspection\_description} package;
  \item complete rover Xacro using primitive shapes;
  \item display launch file and RViz2 configuration;
  \item verified frame tree and screenshots from multiple viewpoints;
  \item dimensional and inertia sanity-check notes.
\end{itemize}

\keypoint{A robot model is an engineering coordinate-system contract. A model
that looks correct from one camera angle can still have a wrong joint axis,
frame origin, collision envelope, or inertia tensor. Verify structure, not only
appearance.}

\subsection{Verification gate}

The Xacro must expand without errors, RViz2 must show all expected links, wheel
motion must occur around the intended axes, and the generated frame tree must
match the documented parent--child relationships. The student must explain why
\texttt{base\_footprint} and \texttt{base\_link} are separate.

\subsection{Mechanical-engineering connection}

This module directly uses CAD-style coordinate reasoning, mass properties,
clearance thinking, and simplified modeling. The student should treat every
dimension and inertia as a justified approximation, not decoration.

\subsection{Suggested commit}

\texttt{feat: add rover URDF and TF model}

\section{Module 3: Gazebo Harmonic simulation and sensor integration}

\textbf{Estimated time:} 4--5 hours

\subsection{Purpose}

Module 3 introduces a computationally light industrial world, spawns the Xacro
rover in \gazebo{}, and connects Gazebo Transport to ROS~2 through
\texttt{ros\_gz}. The rover receives velocity commands and publishes odometry,
joint states, transforms, and a simulated 2D laser scan.

\subsection{Why it follows Module 2}

Simulation adds contacts, friction, gravity, update rates, plugins, and bridge
configuration. These are easier to diagnose after the static model and TF tree
have already been accepted. The student can compare Gazebo behavior against a
known-good RViz representation.

\subsection{Learning objectives}

The student will be able to:

\begin{itemize}
  \item distinguish Gazebo Transport topics from ROS~2 topics;
  \item explain why a bridge is required and verify bridge direction and type;
  \item configure differential-drive wheel joints, separation, radius, and
        odometry publication;
  \item configure a light ray/lidar sensor and connect its frame correctly;
  \item create an SDF world with floor, walls, a narrow passage, two inspection
        stations, a start area, and a marked hazard zone;
  \item command the rover with \texttt{geometry\_msgs/msg/Twist};
  \item inspect odometry signs, frame identifiers, scan ranges, and update rates.
\end{itemize}

\subsection{Core concepts}

SDF worlds, simulation systems/plugins, differential-drive kinematics, wheel
parameters, contact/friction approximations, simulation time, ROS--Gazebo
bridges, odometry, laser scans, and joint-state publication.

\subsection{Principal deliverables}

\begin{itemize}
  \item \texttt{hazard\_inspection\_sim} package and industrial world;
  \item Gazebo-specific Xacro additions;
  \item robot spawning and simulation launch support;
  \item bridge configuration for \texttt{/cmd\_vel}, \texttt{/odom},
        \texttt{/scan}, and \texttt{/joint\_states};
  \item evidence of forward, reverse, and rotation commands;
  \item RViz2 visualization of odometry and laser data.
\end{itemize}

\keypoint{A simulator is another subsystem with its own interfaces and
assumptions. Verify commands, measurements, frames, units, and time at the
boundary instead of treating the simulation as a black box.}

\subsection{Verification gate}

The rover must spawn at the start area, remain physically stable, respond with
the correct sign to manual velocity commands, report plausible odometry, and
show laser returns from simple barriers. The student must compare wheel radius
and separation in the model with the differential-drive configuration.

\subsection{Mechanical-engineering connection}

This module applies idealization and model validation: simplified geometry and
physics are acceptable only when their limitations are stated. The goal is
repeatable software integration, not a high-fidelity tire or contact model.

\subsection{Suggested commit}

\texttt{feat: integrate rover with Gazebo Harmonic}

\section{Module 4: Hazard monitoring, parameters, services, and mission states}

\textbf{Estimated time:} 3--5 hours

\subsection{Purpose}

Module 4 implements a mathematical environmental hazard field based on rover
position. It also establishes mission command services and an explicit
finite-state model before closed-loop waypoint motion is added.

The nominal hazard function is

\begin{equation}
H(x,y) = H_{0} + H_{p}
\exp\left(-\frac{(x-x_{h})^{2} + (y-y_{h})^{2}}{s}\right) + n,
\end{equation}

where $H_{0}$ is the baseline, $H_{p}$ is the peak contribution,
$(x_{h},y_{h})$ is the hazard center, $s$ controls spatial spread, and $n$ is
bounded pseudorandom noise. A configurable seed makes test runs reproducible.

\subsection{Why it follows Module 3}

The hazard node needs a trusted rover position from odometry. Once the odometry
contract is verified, the hazard calculation can be tested independently from
the controller. The finite-state machine is introduced here so that state
transitions can be reasoned about before motion adds timing and geometry.

\subsection{Learning objectives}

The student will be able to:

\begin{itemize}
  \item convert odometry position into a deterministic simulated measurement;
  \item use parameters for hazard location, peak, spread, noise, seed, and rate;
  \item validate parameter ranges and reject invalid sensor inputs;
  \item explain reproducibility and bounded noise in a simulation test;
  \item expose start, reset, and abort operations with standard Trigger services;
  \item represent mission behavior with explicit states and permitted
        transitions;
  \item distinguish a measurement threshold from a complete safety case.
\end{itemize}

\subsection{Core concepts}

Parameter declaration and validation, deterministic random number generation,
services, request/response semantics, custom messages, finite-state machines,
threshold behavior, invalid values, and separation of sensing from decisions.

\subsection{Principal deliverables}

\begin{itemize}
  \item \texttt{hazard\_sensor\_node};
  \item configurable and documented hazard calculation;
  \item \texttt{MissionState.msg} interface;
  \item mission manager with start, reset, and abort services;
  \item explicit states: \texttt{IDLE}, \texttt{MOVING\_TO\_WAYPOINT},
        \texttt{INSPECTING}, \texttt{HAZARD\_DETECTED},
        \texttt{RETREATING}, \texttt{MISSION\_COMPLETE}, and
        \texttt{ABORTED};
  \item command-line tests using synthetic odometry and service calls.
\end{itemize}

\keypoint{Safety-related behavior should be explicit, observable, and
testable. A threshold is only one input to a state decision; stale data, NaN
values, missing updates, reset rules, and service semantics must also be
defined.}

\subsection{Verification gate}

For identical pose, parameters, and random seed, the hazard sequence must be
repeatable. Moving synthetic odometry toward the hazard center must increase
the nominal value. Invalid odometry must not produce a plausible-looking
measurement. Start, abort, and reset calls must produce only documented state
transitions.

\subsection{Mechanical-engineering connection}

The mathematical field resembles a simplified exposure or temperature model.
Parameter sweeps, noise bounds, threshold selection, sensor validity, and
failure-mode reasoning connect directly to test engineering and hazard analysis.

\subsection{Suggested commits}

\begin{itemize}
  \item \texttt{feat: add simulated hazard sensor}
  \item \texttt{feat: add mission services and state model}
\end{itemize}

\section{Module 5: Waypoint control, defensive behavior, and integration}

\textbf{Estimated time:} 4--5 hours

\subsection{Purpose}

Module 5 completes the first integrated mission. A proportional controller
drives toward predefined 2D waypoints, pauses for inspection, and reacts to a
hazard by immediately commanding zero velocity before retreating toward the
previous safe waypoint or start. A logger records synchronized mission
observations, and a bringup package starts the complete demonstration.

For rover pose $(x,y,\theta)$ and waypoint $(x_{w},y_{w})$, the controller uses

\begin{align}
d &= \sqrt{(x_{w}-x)^{2} + (y_{w}-y)^{2}}, \\
\theta_{d} &= \operatorname{atan2}(y_{w}-y, x_{w}-x), \\
e_{\theta} &= \operatorname{normalize}(\theta_{d}-\theta).
\end{align}

Angular velocity is proportional to heading error and linear velocity is
proportional to distance only when alignment is adequate. Both commands are
clamped. This is an educational controller for system integration, not a
production navigation algorithm.

\subsection{Why it follows Module 4}

The controller combines nearly every earlier interface. Implementing it after
the model, simulator, odometry, hazard sensor, services, and states are verified
reduces the number of simultaneous unknowns. Safety transitions already have
defined semantics before they command real simulated motion.

\subsection{Learning objectives}

The student will be able to:

\begin{itemize}
  \item calculate distance, desired heading, and normalized heading error;
  \item implement rotate-then-advance proportional waypoint behavior;
  \item reduce linear speed near a target and saturate both velocity components;
  \item stop on stale odometry, stale hazard data, excessive hazard, or invalid
        numeric values;
  \item connect finite-state decisions to explicit zero, forward, rotational,
        and retreat commands;
  \item log time, pose, yaw, commands, hazard, and mission state safely to CSV;
  \item integrate Gazebo, bridges, nodes, RViz2, parameters, and spawning in one
        launch command;
  \item explain why the waypoint controller is not obstacle-aware navigation.
\end{itemize}

\subsection{Core concepts}

Planar pose extraction, angle normalization, proportional feedback, velocity
saturation, waypoint tolerance, inspection dwell time, watchdog timeouts,
fail-safe stops, path visualization, CSV resource management, launch arguments,
and integration configuration.

\subsection{Principal deliverables}

\begin{itemize}
  \item \texttt{mission\_controller\_node};
  \item \texttt{mission\_logger\_node};
  \item waypoint and safety parameters;
  \item published mission state and planned/visited path;
  \item complete \texttt{hazard\_inspection\_bringup} package;
  \item \texttt{inspection\_demo.launch.py} with RViz and Gazebo GUI arguments;
  \item timestamped CSV output and a documented example.
\end{itemize}

\keypoint{The essential control lesson is not merely ``make the rover move.''
It is ``make every commanded motion conditional on valid, recent inputs and an
explicit mission state, with observable limits and a deterministic stop path.''}

\subsection{Verification gate}

One launch command must start the complete system. The rover must visit the
safe inspection location, stop for the configured duration, encounter rising
hazard on the dangerous leg, publish a zero command, enter the hazard and
retreat states, and stop at a safe destination. RViz2 must show the intended
model, TF, scan, path, and status information. The CSV must contain the required
columns and close cleanly after shutdown.

\subsection{Mechanical-engineering connection}

This module combines closed-loop behavior with limits, tolerances, sampling,
and fault responses. The CSV log plays the role of test instrumentation: it
supports comparison between expected state transitions and observed motion.

\subsection{Suggested commits}

\begin{itemize}
  \item \texttt{feat: implement waypoint mission controller}
  \item \texttt{feat: add hazard-triggered safety behavior}
  \item \texttt{feat: add CSV mission logging}
\end{itemize}

\section{Module 6: Testing, validation, and portfolio communication}

\textbf{Estimated time:} 3--5 hours

\subsection{Purpose}

Module 6 converts a working demonstration into a defensible engineering
project. Pure calculations and state rules receive focused unit tests. ROS~2
CLI procedures exercise the integrated system. A requirements-based validation
matrix records expected and actual results. Documentation, screenshots, a short
video, resume bullets, and statement-of-purpose language are based only on the
verified scope.

\subsection{Why it comes last}

Detailed validation requires stable interfaces and behavior. Portfolio claims
must follow evidence, not precede it. Testing is discussed throughout the
course, but this final module consolidates regression tests and acceptance
records after the first complete mission is functioning.

This does not mean ``testing happens only at the end.'' Every prior module has
a verification gate. Module 6 formalizes those checks, adds repeatable tests,
and reviews the entire repository for consistency.

\subsection{Learning objectives}

The student will be able to:

\begin{itemize}
  \item separate pure mathematics and state rules from ROS node plumbing so
        they can be unit-tested quickly;
  \item write pytest cases for hazard calculation, angle normalization,
        waypoint geometry, hazard transitions, saturation, and invalid data;
  \item perform integration checks using nodes, topics, services, parameters,
        TF, and recorded output;
  \item create a requirement-to-test validation matrix with actual-result and
        pass/fail fields;
  \item perform a consistency review of package, executable, entry-point,
        topic, frame, launch, parameter, dependency, build, run, and test names;
  \item design a factual 60--90 second demonstration;
  \item explain limitations and future extensions without implying that they
        are already implemented.
\end{itemize}

\subsection{Core concepts}

Unit tests, integration tests, boundary cases, fault injection, deterministic
fixtures, acceptance criteria, traceability, regression protection,
reproducibility, technical documentation, and responsible portfolio claims.

\subsection{Principal deliverables}

\begin{itemize}
  \item pytest suite covering the required calculations and failure cases;
  \item integration-test command sequence;
  \item completed or student-fillable validation matrix;
  \item professional root README and six complete guided lessons;
  \item architecture, validation, and learning review documents;
  \item screenshots and a technical 60--90 second demonstration storyboard;
  \item final dependency and naming consistency review.
\end{itemize}

\keypoint{A portfolio project is credible when another person can reproduce
the setup, understand the interfaces, see the limitations, run the tests, and
connect every claim to observable evidence.}

\subsection{Verification gate}

All unit tests must pass, the documented integration scenario must be
repeatable, validation entries must contain actual evidence, and every command
in the README must match an installed package or executable. The student must
be able to explain at least one design tradeoff, one failure mode, one debugging
example, and one limitation without relying on scripted narration.

\subsection{Mechanical-engineering connection}

This module closely parallels verification and validation practice: define a
requirement, select a method, state an expected result, collect an actual
result, assess pass/fail, and retain evidence. Failure-analysis experience is
particularly valuable when discussing stale inputs, invalid values, bridge
misconfiguration, and unsafe control assumptions.

\subsection{Suggested commits}

\begin{itemize}
  \item \texttt{test: add controller and sensor unit tests}
  \item \texttt{docs: complete project documentation}
\end{itemize}

\section{Cross-module learning priorities}

\subsection{Interface literacy}

The student should be able to identify every important node, topic, service,
message type, parameter, and frame, and explain which component owns it. This
is more important than memorizing ROS commands.

\subsection{Coordinate and unit discipline}

Frame identity, SI units, joint axes, pose conventions, and timestamps must be
treated as engineering requirements. Many apparent controller failures are
actually frame, sign, or unit failures.

\subsection{Observability before complexity}

Every subsystem should expose enough information to diagnose it independently.
The course therefore favors explicit mission states, command topics, parameter
files, logs, and simple models over hidden or overly abstract behavior.

\subsection{Defensive behavior}

Stopping on stale, missing, NaN, or threshold-exceeding data is part of the
controller design. The simulation demonstrates a software pattern; it does not
constitute a certified safety mechanism.

\subsection{Technical honesty}

The completed project demonstrates ROS~2 integration, simulation, simple
feedback control, state-based behavior, logging, and validation. Optional later
work may add Nav2, SLAM, more realistic sensor models, automated launch testing,
or physical hardware. Those extensions must not be described as completed
until they are implemented and tested.

\section{Assessment and student participation}

Each module uses four forms of evidence:

\begin{enumerate}
  \item \textbf{Prediction:} state the expected behavior before running.
  \item \textbf{Observation:} collect terminal, RViz2, Gazebo, test, or CSV
        evidence.
  \item \textbf{Explanation:} relate the observation to interfaces, frames,
        parameters, or state logic.
  \item \textbf{Modification:} make a bounded change and predict its effect.
\end{enumerate}

The instructor should not advance automatically. The student reports results
and errors after each module. When an error occurs, the terminal output is used
to identify the smallest justified correction. Replacing a complete package is
not an acceptable substitute for understanding a local fault.

\section{Expected final capability}

After all six modules, the student should be able to launch the project with

\begin{verbatim}
ros2 launch hazard_inspection_bringup inspection_demo.launch.py
\end{verbatim}

and explain the complete chain from Gazebo pose and lidar generation, through
bridged ROS~2 messages and mathematical hazard sensing, to mission-state
decisions, saturated velocity commands, retreat behavior, RViz2 displays, CSV
records, and validation tests.

The strongest portfolio outcome is not the visual demonstration by itself. It
is the student's ability to discuss why the architecture was separated, how
the controller works, which assumptions limit it, how invalid data are handled,
what evidence was collected, and what would have to change for more realistic
navigation or hardware deployment.

\section{Recommended discussion questions}

At the end of the course, the student should be prepared to answer:

\begin{enumerate}
  \item Why were five packages used instead of one package?
  \item How does ROS~2 discover a Python executable and its package resources?
  \item What is the difference between \texttt{base\_link},
        \texttt{base\_footprint}, and \texttt{odom}?
  \item Which data cross the Gazebo--ROS bridge, and in which direction?
  \item Why is a fixed random seed useful when validating the hazard sensor?
  \item How are heading error and velocity saturation calculated?
  \item Which conditions force a stop, retreat, or abort?
  \item What does the CSV demonstrate, and what can it not prove?
  \item Why is this controller not equivalent to Nav2 or autonomous navigation?
  \item Which test or debugging result most strongly changed the design?
\end{enumerate}

\section{Document build command}

This file may be compiled on Ubuntu after installing a standard LaTeX toolchain:

\begin{verbatim}
sudo apt install -y latexmk texlive-latex-base \
  texlive-latex-recommended texlive-fonts-recommended
cd hazard_inspection_ws/docs
latexmk -pdf course_outline.tex
\end{verbatim}

Run \texttt{latexmk -c} afterward if intermediate LaTeX files should be
removed. Compiling this outline is optional and does not affect the ROS~2
workspace build.

\end{document}
